%% kapa: exam=1 round=33 year=2022 period=2 form=A subject=accounting subject_ko=회계학 seq=34 number=74
%% source: Q-Net 감정평가사 기출문제 (cst00309) artlSeq=5209997 file=raw/1008/5209997/2022년도 제33회 감정평가사 1차시험 2교시(감정평가 관계법규, 회계학).pdf
%% conversion: kordoc 4.24.51 -> tools/gen1.py (auto); review=proofread-sonnet; flags=none
\begin{question}{74}
(주)감평은 두 개의 제조부문 X, Y와 두 개의 보조부문 S\textsubscript{1}, S\textsubscript{2}를 운영하고 있으며, 배부 전 부문발생원가는 다음과 같다.
\par{\small\centering\begin{tabular}{|c|c|r|}\hline
\multicolumn{2}{|c|}{부 문} & 부문발생원가 \\ \hline
\multirow{2}{*}{보조부문} & S\textsubscript{1} & ￦90 \\ \cline{2-3}
 & S\textsubscript{2} & 180 \\ \hline
\multirow{2}{*}{제조부문} & X & 158 \\ \cline{2-3}
 & Y & 252 \\ \hline
\end{tabular}\par}
보조부문 S\textsubscript{1}은 보조부문 S\textsubscript{2}에 0.5, 제조부문 X에 0.3, 보조부문 S\textsubscript{2}는 보조부문 S\textsubscript{1}에 0.2의 용역을 제공하고 있다. 보조부문의 원가를 상호배분법에 의해 제조부문에 배부한 후 제조부문 X의 원가가 ￦275인 경우, 보조부문 S\textsubscript{2}가 제조부문 X에 제공한 용역제공비율은?
\begin{choices}
  \item 0.2
  \item 0.3
  \item 0.4
  \item 0.5
  \item 0.6
\end{choices}
\end{question}
\answer{2}
